\section{Introduction}
\label{sec:introduction}

Lossless compression is fundamentally a two-part problem: the encoder must first model the input so that probable symbols receive high estimated probability, and only then can the coder map those probabilities into short bit representations. In practice, the quality of a compressor depends on how well the model captures structure in the data and how efficiently the coding stage turns that structure into a compact stream.

This project was developed in the context of the Algorithmic Information Theory course and evolved through several compression strategies. The repository now contains three especially relevant variants. \texttt{g07-v5} is the clearest balanced milestone and still offers the best overall speed/ratio balance; \texttt{g07-v9} extends the same statistical family and achieves the best compression ratio among the project versions; and \texttt{g07-v8} explores the opposite design point, replacing the statistical pipeline with a very fast LZ77-style matcher.

The goal of this report is therefore not to narrate every version in detail, but to explain the main design choices, summarize the implementation strategy, and present the measured trade-offs on the benchmark corpus. In particular, we use \texttt{g07-v5} as the main case study for the project's best overall compromise, \texttt{g07-v9} as the ratio leader in the current repository snapshot, and \texttt{g07-v8} as the speed-oriented contrast.
